\section{Simple algebras and extended traces}
\label{ext:sec:preliminaries}

The next three sections prove full almost divisibility from
extended-functional comparison. Two consequences of simplicity supply
the error estimates: there are arbitrarily small nonzero Cuntz classes,
and one fixed positive cut has finite positive value at every nonzero
functional that is finite somewhere. The first provides small hereditary
supports; the second measures the errors on all of them. No nuclear
approximation is used.

We work in $B=A\otimes\K$ and identify $\Cu(B)$ with $\Cu(A)$.
Matrix algebras over $B$ are placed in stable corners. These placements
preserve Cuntz classes: if $V$ is the corresponding multiplier isometry,
then $b$ and $VbV^*$ are the two squared moduli of $Vb^{1/2}$.
The trace and rank conventions are those of
Section~\ref{sec:preliminaries}.

\begin{lemma}[Small nonzero classes]
\label{ext:lem:small}
If $A$ is simple and non-elementary, then $B=A\otimes\K$ is simple
and has no minimal projections. For every nonzero $e\in B_+$ and
integer $M\ge1$, there is a nonzero $h\in B_+$ with
\[
 M[h]\le[e].
\]
\end{lemma}

\begin{proof}
Matrix-unit multipliers extract a nonzero matrix entry from any
nonzero ideal of $B$ and put it into every matrix position.
Simplicity of $A$ then implies simplicity of $B$.
If $p\in B$ were minimal, $Bp$ would be a Hilbert space with
inner product $x^*y=\langle x,y\rangle p$. Left multiplication
is faithful by simplicity, and the elements $xy^*$ act as rank-one
operators. Since $p$ generates $B$ as an ideal, they span a dense
subspace of $B$, giving $B\cong\K(Bp)$. The multiplier corner
$1\otimes e_{11}$ would consequently identify $A$ with compact
operators on a Hilbert space, a contradiction.

There are $M$ nonzero mutually orthogonal positives
$e_1,\ldots,e_M$ in $\her(e)$. Indeed, split a nonzero positive
element at two distinct nonzero spectral points. If no such points
exist, the element is a multiple of a projection, whose corner is
not one-dimensional. A shifted nonscalar selfadjoint in that corner
has two distinct positive spectral points and can be split instead.
Recursion gives the required family while preserving orthogonality.

For nonzero $f,g\in B_+$, simplicity gives $fBg\ne0$.
If $0\ne v\in fBg$, then $[v^*v]=[vv^*]$ is a nonzero class
below both $[f]$ and $[g]$. Iterating this observation gives
$h\ne0$ with $[h]\le[e_i]$ for every $i$. Therefore
$M[h]\le\sum_i[e_i]\le[e]$.
\end{proof}

\begin{lemma}[Functional trichotomy]
\label{ext:lem:trichotomy}
Let $A$ be simple and nonzero. Every functional on $\Cu(A)$ is
either zero, infinite on every nonzero class, or positive on every
nonzero class and finite on every class $[(c-\delta)_+]$ with
$c\in(A\otimes\K)_+$ and $\delta>0$.
The third type consists exactly of the functionals finite and
positive at some class.
\end{lemma}

\begin{proof}
Every nonzero $b\in B_+$ is full. Lemma~\ref{lem:pre:cut-placement}
therefore bounds each positive cut of any $c\in B_+$ by finitely
many copies of $[b]$. If $\varphi([b])=0$, preservation of
increasing suprema gives $\varphi([c])=0$ for every $c$.
If $0<\varphi([b])<\infty$, the same cut comparison gives
finiteness on every cut, and the first observation gives positivity
on every nonzero class. These alternatives exhaust all functionals.
\end{proof}

Write $\Fin(A)$ for the third type in
Lemma~\ref{ext:lem:trichotomy}. This notation does not assert finiteness
on all classes. If $\|a\|=1$ and $e=(a-1/2)_+$, then
\begin{equation}
\label{ext:eq:controller}
 0<\varphi([e])<\infty,\qquad
 \varphi([e])\le\varphi([a])
 \quad(\varphi\in\Fin(A)).
\end{equation}

